\section{Introduction}\label{sec:introduction}

A recovery graph specifies which coordinates may determine another
coordinate. It does not specify how a code is represented or computed.
This distinction matters in the coding approach to matrix transposition:
a structural property of a machine's dependency graph becomes useful
only when a coding theorem applies to its actual code fibres.
The companion paper~\cite{BrendeDirected} obtains long matched-port
distances from native tape geometry and studies several coding bounds
with additional premises. Here we examine two such premises more closely:
algebraic structure in regular covers, and local generation of a linear code.

The cover results combine two settings. For a fixed template repeated
over finite group quotients, homogeneous recovery has an asymptotic
half-rate theorem uniform over coefficient fields and values at fixed
width. We identify what its filtration argument needs when width grows.
For the prescribed Clebsch adjacency check, exact finite bounds instead
control binary towers and specified regular covers, with sharpness and
rigidity and without a separating-chain hypothesis. For local generation
the graph need not be a cover at all. Locally supported codewords
are forced into the image of a check operator, while the whole code lies
in its kernel. Rank-nullity gives a quantitative restriction on their span.

\subsection{The results}
The principal conclusions are the following.
\begin{enumerate}[label=\textup{(\arabic*)}]
\item \textbf{Sharp finite Clebsch-cover bounds.}
Every connected binary-lift tower above Clebsch has adjacency-code
dimension at most $N/2+2$. Either equality and signing rigidity persist,
or the dimension drops to at most half and stays there. The constant
two is sharp at every level. We extend the bound to regular deck groups
$P\times D$, for a $2$-group $P$ and odd abelian $D$, and to
$P\times S_3$; every connected three-sheet cover has dimension at most
$26$. A named rigid three-sheet seed supports further binary towers.
Reciprocal blocks have rank at least $6r$, with equality exactly for
flat gains and no rank $6r+1$. A concrete counterexample limits a
stronger proposed extension
(\cref{sec:clebsch-finite,sec:clebsch-block-proof,sec:clebsch-small}).
\item \textbf{Uniform half rate on fixed covers.}
Under the precise ambient-group and chain hypotheses in
\cref{sec:model}, paired-port disconnection in the infinite incidence
cover forces asymptotic rate at most $1/2$ for homogeneous linear checks
of fixed total width. One threshold works for every coefficient tuple
over every field of the fixed characteristic
(\cref{thm:fixed-width}). The coefficients and fields may therefore vary
along the sequence.
\item \textbf{An obstruction to excess rate at growing width.}
Let $\alpha_i(k)$ be the fraction of internal dimension in simple
coefficient-module factors larger than $k$, after algebraic closure.
For each fixed $k$ and $\varepsilon>0$, eventually
\[
 r_i\le(1-\alpha_i(k))(1/2+\varepsilon)+\alpha_i(k).
\]
Thus a sustained excess $\delta$ over half rate requires at least
$2\delta$ of the dimension in factors exceeding every fixed cutoff
(\cref{thm:factors,cor:factor-excess}).
\item \textbf{Uniform zero rate for a stricter check class.}
On the independent-chord covers of a fixed graph with a spanning
$2$-factor, every homogeneous scalar check with nonzero edge weights
has rate tending to zero. The convergence is uniform in the field,
weights, and diagonal entries. It also holds for arbitrary-width
blocks with a common full flag and invertible edge maps
(\cref{thm:weighted-zero,cor:flag-zero}).
\item \textbf{The cost of removing sheet dependence.}
Any sheet-dependent linear check has an exact homogeneous inflation
with the same rate. Its internal width grows by the sheet count, and
its coefficient algebra can be a full matrix algebra of that growing
size. Hence this conversion does not supply the small-factor premise
(\cref{prop:inflation,prop:large-factor}).
\item \textbf{A local-generation deficit.}
In characteristic two and with no odd cycle of length at most $2R+1$,
\[
 \dim\mathcal C+\dim\mathcal L_R(\mathcal C)
       \le\sum_v\dim\pi_v(\mathcal C).
\]
An above-half code must therefore have substantial dimension outside
the span of its locally supported words (\cref{thm:local-generation}).
If a nonlinear Cartesian-source encoder has a recoverable affine hull,
nonlocal sources must pay for excess support size or entropy
(\cref{thm:local-sources}). All-local sources generate a hull of at
most half rate.
\item \textbf{Group-word extensions of the cover bounds.}
Ordinary-word recovery over finite $p$-groups reduces to the scalar
exponent check modulo $p$. Homogeneity of that check transfers the
fixed-width half bound, and nonzero homogeneous edge exponents transfer
the independent-chord zero-rate bound. The alphabet groups may grow
without affecting the threshold (\cref{cor:group-cover}).
\end{enumerate}

\subsection{Relation to existing work}
Storage codes and index coding are established subjects. In particular,
triangle-free graphs can support binary linear codes with rate approaching
one~\cite{BargZemor,BargSchwartzYohananov,Zabokritskiy}.
We make no general half-rate claim for triangle-free graphs. Our local
theorem instead constrains the span of local codewords on such graphs;
recovery locality and generator locality are different conditions.

The cover argument uses standard free-group subgroup separability,
presented here by finite Schreier-graph completion; see
Stallings~\cite{Stallings} for the classical graph method. The
free-group algebra input is Cohn's Sylvester-domain theorem
\cite{Cohn}, as recorded in~\cite[Introduction]{JaikinSouza}.
The asymptotic input is precisely Jaikin-Zapirain and Souza's
Theorem~8.1~\cite{JaikinSouza}. We apply that theorem to the specified
quotients, then prove a separate uniformization over coefficient fields
and parameters. We do not claim a new general rank-approximation theorem.
The parameter argument uses the finite-residue-field consequence of the
weak Nullstellensatz~\cite{Stacks00FV}.

The first paper formulates the tape problem and its coding program
\cite{BrendeTapeCoding}; the subsequent local-barrier paper separates
cheap encoding from cheap local recovery~\cite{BrendeLocalCoding}.
The fourth paper supplies depth reduction and the reduction from a
correct trace to a storage-code fibre~\cite{BrendeCurved}.
The directed-recovery companion proves native paired-distance bounds
and an all-width matrix-label criterion on a different family
\cite{BrendeDirected}. Our cover theorems neither subsume that
all-width criterion nor establish its hypotheses for arbitrary cones.
They remove field and coefficient restrictions in a specified cover
model and expose the remaining width obstruction.

\paragraph{Proof status and attribution.}
Exact homogenization and the dimension interfaces, including the quotient
deficit and split-generator criterion, are Lean-verified. Their formal
statements are authoritative; the English outlines are guides to them.
The finite and asymptotic cover results, graph-local argument, and source
and group-word applications have complete written proofs but are not formalized.
The named rigid seed includes an explicit binary matrix-inverse certificate
with a portable exact checker; that checker is not a Lean proof.
In particular, the split-generator
interface does not itself verify the graph parity and boundary argument.
Every numbered result states its status. We distinguish the applications
developed here from their standard algebraic inputs; external priority for
the resulting combinations has not been established.

\subsection{Connection to transposition}
Keep the series' notation: $d$ is the matrix side length,
$b(d)=\max\{1,\lceil\log_2d\rceil\}$, and $L=d^2b(d)$.
Statement~(T) asks for one fixed multitape machine transposing these
inputs in $O(L)$ time. A coding lower-bound route needs a code deficit
that exceeds a paid genie cost. The present theorems can supply such a
deficit only after their representation and linearity hypotheses have
been proved for the relevant fibres. They are also tests for constructions:
an above-half candidate in these regimes must escape the stated module
or local-generation restrictions. The missing native implications are
listed in \cref{sec:roadmap}; (T) remains open.
